\makeabstract
{
Biochemical Characterization of Human LIG1 Syndrome Variants as Possible Players in Cancer Expression 
}
{
Kellen Garris (1), David Beier (2), \& Patrick O’Brien (2)
}
{
\textsuperscript{1}Department of Neuroscience, Amherst College, Amherst, MA, 0100\\
\textsuperscript{2}University of Michigan Department of Biological Chemistry, Ann Arbor, MI 48109 USA
}
{
Human DNA ligase 1 (LIG1) uses ATP and magnesium in the final steps of nucleotide excision repair, long-patch base excision repair, and normal DNA replication. Mutation of this enzyme can lead to human LIG1 syndrome, characterized by immunodeficiency and developmental delays. There is growing evidence for the involvement of R305Q, R771G, and R771W LIG1 syndrome mutations in genome instability and non-specific patient cancers. Determining whether rare LIG1 mutants increase cell mutagenesis by shifting enzyme fidelity would underscore LIG1’s involvement in cancer formation and highlight a relatively new therapeutic target for cancer prevention. 
}
{
Steady-state kinetics competition assays were run before using urea polyacrylamide gel electrophoresis (PAGE) for product separation and quantification. Assays were run at 37℃ and quenched in FLS. An analysis of initial turnover rates allowed me to look at how well these different LIG1 mutants discriminated against sealing oxidized (damaged) DNA single-stranded breaks and preferentially closed the backbone of healthy canonical DNA. 
}
{
All mutants were able to better discriminate against G:A mismatches compared to WT LIG1. As cofactor increased, fidelity decreased for all LIG1 species except R771G. 
}
{
The LIG1 variants tested all have higher fidelity than WT using a 3’ 8oxoG:A pair. These mutants appear unlikely to increase cancer risk than WT LIG1. While R771W and R305Q are sensitive to increasing free [Mg2+], R771G appears resistant to increasing free cofactor. This outcome favors a LIG1 syndrome model in which longer-living single-stranded breaks are a source of mutations. 
}

\newpage
